\appendix
\newcommand{\scenariodiagram}[4]{
\begin{figure}[H]
	\begin{center}
		\IfFileExists{#1}{
			\includegraphics[width=0.97\textwidth,height=0.78\textheight,keepaspectratio]{#1}
		}{
			\fbox{\parbox[c][0.58\textheight][c]{0.90\textwidth}{\centering #4}}
		}
		\caption{#2}
		\label{#3}
	\end{center}
\end{figure}
}

\setcounter{figure}{0}
\renewcommand{\thefigure}{А\arabic{figure}}

\specialsection{Приложение А}

В приложении А приведены диаграммы последовательности ключевых сценариев системы.

\scenariodiagram{diagrams/sequence-happy-path.png}{Диаграмма последовательности основного сценария: создание задачи и фиксация связанного контекста}{fig:appendix-sequence-happy-path}{Диаграмма последовательности основного пользовательского сценария.}

\clearpage

\scenariodiagram{diagrams/sequence-access-denied.png}{Диаграмма последовательности альтернативного сценария: отказ в изменении статуса задачи агентом}{fig:appendix-sequence-access-denied}{Диаграмма последовательности сценария обработки ошибки авторизации агентного запроса.}

\clearpage

\scenariodiagram{diagrams/sequence-integration-sync.png}{Диаграмма последовательности сложного сценария: синхронизация артефактов GitHub и файлового хранилища}{fig:appendix-sequence-integration-sync}{Диаграмма последовательности интеграционного сценария с внешними системами.}

\clearpage

\setcounter{figure}{0}
\renewcommand{\thefigure}{Б\arabic{figure}}

\specialsection{Приложение Б}

В приложении Б приведена BPMN-диаграмма ключевого бизнес-процесса системы.

\scenariodiagram{diagrams/bpmn-task-lifecycle.png}{BPMN-диаграмма процесса обработки задачи}{fig:appendix-bpmn-task-lifecycle}{BPMN-диаграмма процесса обработки задачи от постановки до завершения.}

\clearpage
